\begin{table}[tbp]
\centering
\caption{Complete frozen aggregate comparison of all benchmark methods.}
\label{tab:all-methods}
\scriptsize
\resizebox{\linewidth}{!}{%
\begin{tabular}{lrrrr}
\toprule
Method & Local-effect MAE & AUPRC & Macro-F1 & Regional FPR \\
\midrule
Standard SC & 0.09983 & 0.91476 & 0.81641 & 0.140 \\
MetaShift fixed & 0.10344 & 0.89832 & 0.79515 & 0.150 \\
MetaShift CV & 0.09742 & 0.90178 & 0.80916 & 0.135 \\
Nearest-neighbor DiD & 0.11150 & 0.88443 & 0.78852 & 0.075 \\
Bayesian mean shift & 0.24672 & 0.50116 & 0.49962 & 0.542 \\
Before-after median & 0.25330 & 0.50125 & 0.50024 & 0.505 \\
CUSUM & N/A & 0.50117 & 0.49720 & 0.581 \\
Rolling MAD & N/A & 0.50112 & 0.49427 & 0.402 \\
PELT & N/A & 0.50025 & 0.49976 & 0.522 \\
\bottomrule
\end{tabular}
}
\par\smallskip
\begin{minipage}{0.94\linewidth}
\footnotesize\textit{Note.} All values are aggregates over the held-out stable-regime evaluation partition. Lower MAE and regional FPR are preferable; higher AUPRC and macro-F1 are preferable. N/A means that the method supplies no local-effect magnitude estimate for that perturbation design, not that the method was omitted.
\end{minipage}
\end{table}
